\section{System Implementation}
\label{sec:system_implementation}

This chapter provides an overview of the implemented system, including the functional requirements addressed in this release, the adopted infrastructure and programming languages, and the organization of the source code.

\subsection{Implemented Requirements}
\label{subsec:implemented_requirements}

The implementation focuses on the core functional requirements defined in the Design Document (DD), prioritizing the "Phase 0" and "Phase 1" goals to establish a viable prototype. Below is a summary of the implemented requirements.

\subsubsection{User Management (F1)}
\begin{itemize}
    \item \textbf{Registration (F1.1):} Users can create accounts using a valid email and password.
    \item \textbf{Login (F1.2):} Secure authentication using JWT tokens and bcrypt password hashing.
    \item \textbf{Profile Management (F1.3):} Users can view and simple update their profile information.
\end{itemize}

\subsubsection{Trip Management (F2)}
\begin{itemize}
    \item \textbf{Manual Trip Recording (F2.1):} Users can start and stop a trip recording manually.
    \item \textbf{Trip Statistics (F2.5):} The system calculates distance and duration for recorded trips.
    \item \textbf{Trip History (F2.7):} Users can retrieve their past trips.
\end{itemize}

\subsubsection{Path Management (F3)}
\begin{itemize}
    \item \textbf{Path Insertion (F3.1):} Users can insert new paths with geospatial data (coordinates) and metadata (name, status).
    \item \textbf{Path Search (F5.1 - F5.2):} Basic search functionality to find paths near a location or within a region.
\end{itemize}

\subsubsection{Automated Detection (F4)}
\begin{itemize}
    \item \textbf{Sensor Data Acquisition:} The mobile app collects data from GPS, Accelerometer, and Gyroscope (F4.1), which is sent to the backend for processing.
    \item \textbf{Obstacle Detection (F4.3):} Basic server-side logic to identify potential anomalies (potholes) from sensor streams.
\end{itemize}

\subsubsection{Exclusions}
Requirements related to **Social Features (F7)** were excluded from this release to focus on the core functionality of data collection and path management. The "Bottom-Up" strategy prioritized establishing robust data layers before implementing higher-level social interactions.

\subsection{Implemented Infrastructure}
\label{subsec:implemented_infrastructure}
To support the system's functional requirements and ensure scalability, the following infrastructure components have been implemented:
\begin{itemize}
    \item \textbf{PostgreSQL:} Serves as the primary relational database management system, storing persistent data such as user profiles, trip histories, and path metadata. PostGIS is enabled for spatial indexing.
    \item \textbf{Redis:} Utilized as a high-performance in-memory caching layer, specifically for managing active user sessions and real-time trip data to reduce database load.
    \item \textbf{RabbitMQ:} Implemented as a message broker to handle asynchronous communication between services, decoupling heavy processing tasks (like path statistics calculation) from the main request flow.
\end{itemize}

\subsection{Programming Languages}
\label{subsec:programming_languages}
The implementation leverages a unified language stack to streamline development:
\begin{itemize}
    \item \textbf{Backend:} Developed in \textbf{Node.js} (JavaScript), utilizing Express.js for the REST API and standard libraries for service integrations.
    \item \textbf{Frontend (Mobile):} Built with \textbf{React Native} (JavaScript) using the Expo framework, allowing for a cross-platform mobile application that interacts with the backend services.
\end{itemize}

\subsection{Source Code Structure}
\label{subsec:source_code_structure}

The project uses a **Monorepo** structure. This organization was chosen to keep backend and frontend code synchronized, share configuration (linting, prettier), and simplify dependency management.

\textbf{Directory Structure:}
\begin{itemize}
    \item \textbf{\texttt{code/backend/}}: Contains the server-side Node.js application.
    \begin{itemize}
        \item \texttt{src/controllers/}: Handles HTTP requests and response formatting.
        \item \texttt{src/services/}: Contains the business logic for Users, Trips, and Paths.
        \item \texttt{src/repositories/}: Encapsulates raw Database/SQL queries.
        \item \texttt{migrations/}: SQL scripts to create and update database schema.
    \end{itemize}

    \item \textbf{\texttt{code/frontend/}}: Contains the React Native mobile application.
    \begin{itemize}
        \item \texttt{src/pages/}: React components representing full application screens.
        \item \texttt{src/components/}: Reusable UI elements (Buttons, Cards).
        \item \texttt{src/services/}: API clients for backend communication.
        \item \texttt{src/sensors/}: Modules interfacing with device hardware (GPS, Gyroscope).
    \end{itemize}

    \item \textbf{\texttt{code/docker-compose.yml}}: Describes the infrastructure services (Postgres, Redis, RabbitMQ) required to run the system locally.
\end{itemize}
